\documentclass{article}
\usepackage{graphicx} % Required for inserting images
\usepackage{amsmath}
\title{Lesson 8 Homework (Extra Challenge) - Proof Writing}
\author{William Wang}
\date{October 2025}

\begin{document}

\maketitle

\section{}
Skipped
\section{}
Since we have $\sec x + \tan x = 22/7$, we also have from the Pythagorean Identity
\[\sec^2 x - \tan^2 x = 1 \implies (\sec x - \tan x)(\sec x + \tan x) = 1\]
\[\implies \sec x + \tan x = 22/7 \text{ and } \sec x - \tan x = 7/22\]
Then we can just solve for \(\sec x\) and $\tan x$
\[2\sec x = \frac{22}{7} + \frac{7}{22} = \frac{533}{154} \implies \sec x = \frac{533}{308}\]
\[2\tan x = \frac{22}{7} - \frac{7}{22} = \frac{435}{154} \implies \tan x = \frac{435}{308}\]
Since $\sec x = \frac{1}{\cos x}$, we can find $\cos x = \frac{308}{533}$ and since $\tan x= \frac{\sin x}{\cos x}$, we have 
\[\sin x = \tan x \cdot \cos x = \frac{435}{533}\]
Then we can solve for $\csc x + \cot x$ since $\csc x = \frac{1}{\sin x}$ and $\cot x = \frac{1}{\tan x}$
\[\csc x + \cot x = \frac{533}{435} + \frac{308}{435} = \frac{841}{435} = \frac{29}{15}\]
\section{}
Skipped
\end{document}
